%! Graphing Radical Functions \& Transformations — Unit 5, Lesson 5.4 — Group Activity (Answer Key)
\documentclass[10pt]{article}
\usepackage{algebra2-article}
\usepackage{algebra2-key}   % pulls in -boxes; do NOT also load -boxes
\newcommand{\TallMath}[1]{$\displaystyle #1\rule[-1.4em]{0pt}{3.2em}$}
% In-formula answer slot (\ans is text-mode and must never appear inside $...$).
\newcommand{\mans}[1]{{\color{keyred}\mathbf{#1}}}
\newcommand{\lgrid}[4]{%
  \draw[step=1, linegray!40, very thin] (#1,#3) grid (#2,#4);
  \draw[->, semithick, charcoal] (#1,0)--(#2,0) node[below right, font=\scriptsize]{$x$};
  \draw[->, semithick, charcoal] (0,#3)--(0,#4) node[left, font=\scriptsize]{$y$};}
% Cube root a*cbrt(x-h)+k: pgfmath has no cbrt, so plot the two branches separately.
% #1=a  #2=h  #3=k  #4=xmin  #5=xmax
\newcommand{\cbrtplot}[5]{%
  \draw[very thick, forest, domain=#2:#5, samples=70, smooth]
    plot (\x,{(#1)*pow((\x)-(#2),1/3)+(#3)});
  \draw[very thick, forest, domain=#4:#2, samples=70, smooth]
    plot (\x,{-(#1)*pow((#2)-(\x),1/3)+(#3)});}

\begin{document}

\pageheader{Unit 5, Lesson 5.4}{Group Activity --- Answer Key}
\namepartnerperiod

\vspace{0.08in}

\begin{headlinebox}{forestbg}
\small\textbf{Follow the Anchor.} Every radical graph is built around one point. Find it, and the
domain, the range, and the direction of the arm all come free. With your partner, work each problem
\emph{and} say out loud which part of the equation produced each answer. Write every domain and
range in \textbf{interval notation}.
\end{headlinebox}

\vspace{0.06in}

\begin{tcolorbox}[colback=white, colframe=black!40, title=\textbf{Tier R --- Remediate},
    fonttitle=\bfseries, arc=1.5mm, left=3mm, right=3mm, top=2mm, bottom=2mm, breakable]
The solid graph is $r(x)=\sqrt{x-2}+1$. The dashed graph is the parent $y=\sqrt{x}$, shown so you
can see what moved.
\begin{center}
\begin{tikzpicture}[scale=0.46]
  \lgrid{-2}{13}{-2}{6}
  \begin{scope}
    \clip (-2,-2) rectangle (13,6);
    \draw[very thick, slate, dashed, domain=0:13, samples=100, smooth] plot (\x,{sqrt(\x)});
    \draw[very thick, forest, domain=2:13, samples=100, smooth] plot (\x,{sqrt((\x)-2)+1});
  \end{scope}
  \fill[slate] (0,0) circle (3.4pt);
  \fill[forest] (2,1) circle (4pt) node[left, font=\scriptsize]{$(2,1)$};
  \fill[charcoal] (3,2) circle (2.6pt);
  \fill[charcoal] (6,3) circle (2.6pt) node[above left, font=\scriptsize]{$(6,3)$};
  \fill[charcoal] (11,4) circle (2.6pt);
\end{tikzpicture}
\end{center}
\begin{enumerate}[leftmargin=1.7em, itemsep=7pt]
  \item Compare $r(x)=\sqrt{x-2}+1$ with the parent $y=\sqrt{x}$. \\
        $h=\mans{2}$ \quad $k=\mans{1}$ \quad so the graph moved right \ans{2}
        and up \ans{1}.
  \item The \textbf{anchor point} (for a square root, the endpoint) is \ans{$(2,1)$}. \\
        Check it against the picture: the dashed graph starts at $(0,0)$ and the solid one starts at
        \ans{$(2,1)$}.
  \item \textbf{Domain:} \ans{$[2,\infty)$} \quad \textbf{Range:} \ans{$[1,\infty)$}
  \item Now find the domain \emph{the Lesson 5.0 way}: set the radicand $\ge 0$ and solve. \\
        $x-2\ge 0$, so $x\ge\mans{2}$. \; Does this match your answer to \#3?
        \ans{yes} \\
        Explain, in one sentence, why these two methods can never disagree. \ans{Shifting the graph
        right $2$ and refusing every input below $2$ are the same event described twice.}
  \item Is the anchor an absolute maximum or an absolute minimum? \ans{absolute minimum} \\
        Which way does the arm leave the anchor --- up or down? \ans{up}
  \item \textbf{End behavior.} As $x\to+\infty$, $r(x)\to\mans{+\infty}$. \\
        As $x\to-\infty$: careful --- say exactly what happens. \ans{Nothing --- the graph stops at
        $x=2$. There is no left arm.}
  \item Does $r$ have a $y$-intercept? \ans{no} \quad Why or why not? \ans{$x=0$ is not in the
        domain; $r(0)=\sqrt{-2}+1$ is not real.}
\end{enumerate}
\end{tcolorbox}

\begin{tcolorbox}[colback=white, colframe=black!40, title=\textbf{Tier A --- Approaching Proficiency},
    fonttitle=\bfseries, arc=1.5mm, left=3mm, right=3mm, top=2mm, bottom=2mm, breakable]
\textbf{Part 1 --- No graphs allowed.} Everything below can be read straight off the equation.
\begin{center}
\renewcommand{\arraystretch}{1.5}
\begin{tabularx}{\linewidth}{@{}>{\bfseries}p{0.20\linewidth} X X X X@{}}
 & $p(x)=\sqrt{x+3}-2$ & $q(x)=-\sqrt{x-1}+3$ & $s(x)=\sqrt[3]{x-2}+1$ & $t(x)=\tfrac12\sqrt{x}-4$ \\
\hline
$a$ & \ans{$1$} & \ans{$-1$} & \ans{$1$} & \ans{$\tfrac12$} \\
$h$ & \ans{$-3$} & \ans{$1$} & \ans{$2$} & \ans{$0$} \\
$k$ & \ans{$-2$} & \ans{$3$} & \ans{$1$} & \ans{$-4$} \\
Anchor point & \ans{$(-3,-2)$} & \ans{$(1,3)$} & \ans{$(2,1)$} & \ans{$(0,-4)$} \\
Domain & \ans{$[-3,\infty)$} & \ans{$[1,\infty)$} & \ans{all reals} & \ans{$[0,\infty)$} \\
Range & \ans{$[-2,\infty)$} & \ans{$(-\infty,3]$} & \ans{all reals} & \ans{$[-4,\infty)$} \\
Anchor is a max, a min, or neither? & \ans{min} & \ans{max} & \ans{neither} & \ans{min} \\
\end{tabularx}
\end{center}

\vspace{4pt}
\textbf{Part 2 --- Write the equation.} The graph below is a square root function in the form
$y=a\sqrt{x-h}+k$.
\begin{center}
\begin{tikzpicture}[scale=0.44]
  \lgrid{-1}{14}{-5}{5}
  \begin{scope}
    \clip (-1,-5) rectangle (14,5);
    \draw[very thick, forest, domain=4:14, samples=100, smooth] plot (\x,{2*sqrt((\x)-4)-3});
  \end{scope}
  \fill[forest] (4,-3) circle (4pt) node[below left, font=\scriptsize]{$(4,-3)$};
  \fill[charcoal] (5,-1) circle (2.6pt);
  \fill[charcoal] (8,1) circle (2.6pt) node[above left, font=\scriptsize]{$(8,1)$};
  \fill[charcoal] (13,3) circle (2.6pt);
\end{tikzpicture}
\end{center}
\begin{enumerate}[label=(\alph*), leftmargin=1.9em, itemsep=5pt]
  \item Anchor point: \ans{$(4,-3)$} \quad so $h=\mans{4}$ and $k=\mans{-3}$.
  \item Substitute the point $(8,1)$ and solve for $a$. \; $a=\mans{2}$

        {\small\color{keyred}$1=a\sqrt{8-4}-3 \;\Rightarrow\; 1=a\sqrt{4}-3 \;\Rightarrow\;
        1=2a-3 \;\Rightarrow\; 4=2a \;\Rightarrow\; a=2$}
  \item \textbf{Equation:} $y=$ \ans{$2\sqrt{x-4}-3$}
  \item Check your equation against the third point $(13,3)$. Show the substitution.
        \ans{$2\sqrt{13-4}-3=2\sqrt{9}-3=2(3)-3=3$ \checkmark}
\end{enumerate}

\vspace{5pt}
\textbf{Part 3 --- Justify.}
\begin{enumerate}[label=\textbf{J\arabic*.}, leftmargin=2.4em, itemsep=6pt]
  \item $p$ and $q$ are both square root functions, but only one of them has an absolute
        \emph{maximum}. Which one, and which single symbol in its equation is responsible?
        \ansline{$q$, and the responsible symbol is the \emph{minus sign} in front of the radical
        ($a=-1$). A negative $a$ makes every output the opposite of the parent's, so the arm leaves
        the anchor going \emph{down} instead of up --- which makes the anchor the highest point on
        the graph rather than the lowest. The $+3$ has nothing to do with it; it only says
        \emph{where} that highest point sits.}
  \item $s(x)=\sqrt[3]{x-2}+1$ has an $h$ and a $k$ just like the others, and its graph is still
        built around the point $(2,1)$. So why is that point \emph{not} an endpoint, and why does
        knowing it tell you nothing about the domain? \ansline{Because the index is \emph{odd}. An
        endpoint exists only when some inputs are forbidden --- the graph has to stop somewhere. A
        cube root accepts every real number, so nothing stops it: the graph runs through $(2,1)$ and
        keeps going in both directions. The point is still the anchor (it is where the curve
        flattens and changes the way it bends --- the \textbf{inflection point}), but the domain is
        all real numbers no matter what $h$ is, so $h$ tells you only \emph{where} the graph is, not
        \emph{where it is allowed to be}.}
\end{enumerate}
\end{tcolorbox}

\begin{tcolorbox}[colback=white, colframe=black!40, title=\textbf{Tier E --- Extension},
    fonttitle=\bfseries, arc=1.5mm, left=3mm, right=3mm, top=2mm, bottom=2mm, breakable]
\textbf{Part 1 --- The hidden $h$.} None of these radicands is in the form $x-h$. Factor out the
coefficient of $x$ first, then read the graph off the result.
\begin{center}
\renewcommand{\arraystretch}{1.6}
\begin{tabularx}{\linewidth}{@{}l X X X@{}}
\hline
\textbf{Function} & \textbf{Factored form} & \textbf{Rewritten as $a\sqrt[n]{x-h}+k$} & \textbf{Anchor \& domain} \\
\hline
$\sqrt{4x}$ & \ans{$\sqrt{4\cdot x}$} & \ans{$2\sqrt{x}$} & \ans{$(0,0)$; $[0,\infty)$} \\
$\sqrt{9x-18}$ & \ans{$\sqrt{9(x-2)}$} & \ans{$3\sqrt{x-2}$} & \ans{$(2,0)$; $[2,\infty)$} \\
$\sqrt[3]{8x+8}$ & \ans{$\sqrt[3]{8(x+1)}$} & \ans{$2\sqrt[3]{x+1}$} & \ans{$(-1,0)$; all reals} \\
\hline
\end{tabularx}
\end{center}
A classmate looks at $\sqrt{9x-18}$ and says ``the graph starts at $x=18$.'' What went wrong, and
what is the real answer? \ansline{They read the constant straight out of the radicand without
factoring, as though the radicand were $x-18$. It is not: the $x$ has a coefficient of $9$. Factoring
gives $\sqrt{9(x-2)}=3\sqrt{x-2}$, so the graph starts at $x=\mathbf{2}$ --- confirmed the old way by
$9x-18\ge0 \Rightarrow x\ge2$.}

\vspace{5pt}
\textbf{Part 2 --- Two names for one graph.} The SOL lists $f(kx)$ (squeeze the input) and $kf(x)$
(stretch the output) as \emph{different} transformations. For radical functions they collapse into
one.
\begin{enumerate}[label=(\alph*), leftmargin=1.9em, itemsep=6pt]
  \item Explain, using the product property $\sqrt[n]{ab}=\sqrt[n]{a}\cdot\sqrt[n]{b}$, why
        $y=\sqrt{kx}$ is the same graph as $y=\sqrt{k}\cdot\sqrt{x}$ for every $k>0$.
        \ansline{The product property splits the radicand: $\sqrt{kx}=\sqrt{k}\cdot\sqrt{x}$. Since
        $k$ is a constant, $\sqrt{k}$ is just a number, so the right side is the parent multiplied
        by that number. The two expressions are equal at every $x\ge0$, so they cannot be two
        graphs --- they are one graph with two descriptions.}
  \item Rewrite the claim with \textbf{rational exponents} (Lesson 5.1) to see it in one line:
        $(kx)^{1/n}=$ \ans{$k^{1/n}$} $\cdot$ \ans{$x^{1/n}$}. \\
        So compressing the input of $\sqrt[n]{x}\,$ by a factor of $k$ is the same as stretching the
        output by a factor of \ans{$\sqrt[n]{k}$}.
  \item Squeezing $y=\sqrt{x}$ horizontally by $16$ gives $y=\sqrt{16x}$. What single vertical
        stretch produces the identical graph? \ans{a stretch by $4$, since $\sqrt{16x}=4\sqrt{x}$}
\end{enumerate}

\vspace{5pt}
\textbf{Part 3 --- Where the trick breaks.} Now try the same move on $y=\sqrt{x}+1$, which is the
parent \emph{plus a vertical shift}.
\begin{center}
\renewcommand{\arraystretch}{1.3}
\begin{tabular}{|l|c|c|c|c|}
\hline
$x$ & $0$ & $1$ & $4$ & $9$ \\ \hline
Squeeze the input: \; $\sqrt{9x}+1$ & \ans{$1$} & \ans{$4$} & \ans{$7$} & \ans{$10$} \\ \hline
Stretch the output: \; $3\left(\sqrt{x}+1\right)$ & \ans{$3$} & \ans{$6$} & \ans{$9$} & \ans{$12$} \\ \hline
\end{tabular}
\end{center}
\begin{enumerate}[label=(\alph*), leftmargin=1.9em, itemsep=6pt, start=4]
  \item Are the two rows the same? \ans{no} \quad By how much do they differ, at every
        single $x$? \ans{by exactly $2$}
  \item Explain what went wrong. Which part of $\sqrt{x}+1$ refuses to be handled by the product
        property, and why? \ansline{The $+1$. The product property splits a \emph{product} inside a
        radical --- it says nothing about a term sitting \emph{outside} one. Squeezing the input
        touches only the radical, giving $3\sqrt{x}+1$, while stretching the output multiplies
        \emph{everything}, giving $3\sqrt{x}+3$. The radical parts agree; the constant does not, and
        the gap is $3(1)-1=2$ at every $x$.}
  \item Write a one-sentence rule a classmate could trust: \emph{for which radical functions can you
        trade a horizontal size change for a vertical one, and for which can you not?}
        \ansline{You can trade them freely for a \emph{pure} radical --- anything of the form
        $a\sqrt[n]{kx}$, with nothing added on --- because the constant walks out from under the
        radical; the moment there is a vertical shift $+k$ (or any term outside the radical), the
        two transformations stop agreeing.}
\end{enumerate}
\end{tcolorbox}

\begin{teachernote}
\small
\textbf{Tier R} exists to make one comparison happen: item~3 and item~4 ask for the same domain by
two different routes. Do not let a group skip item~4 because they ``already have the answer'' ---
answering it twice \emph{is} the point, and the one-sentence explanation is the lesson's thesis in
the students' own words. Item~6 recycles 5.0's most-missed idea (``the graph stops'' is a complete
answer); item~7 catches anyone still hunting for a $y$-intercept outside the domain.

\textbf{Tier A, Part 1} is deliberately arranged so that $t(x)=\tfrac12\sqrt{x}-4$ is the only one
with a fractional $a$ and $s$ is the only cube root --- every column differs from every other in at
least two rows, so a group that pattern-matches instead of reading will produce a visibly wrong
table. The row that separates proficiency from fluency is the last one: \emph{neither} is the
correct answer for $s$, and students reliably want to write ``min.''

\textbf{Part 2 is the SOL's harder direction} (\textbf{A2.F.1b}, write the equation from a graph).
Insist on the order --- anchor first, then $a$ --- and on the check in (d). Groups that hunt for $a$
first end up with two unknowns and stall.

\textbf{Tier E, Part 3 is the item worth circulating for.} Parts 1 and 2 will convince a fast group
that horizontal and vertical size changes are ``always the same thing,'' which is false and would
carry a real misconception into Unit 6. The table forces the counterexample into the open before
anyone can generalize, and the constant gap of $2$ is a satisfying thing to notice. If a group
finishes early, ask them the follow-up: \emph{what would the gap be for $\sqrt{x}+5$?} (It is
$3(5)-5=10$.)
\end{teachernote}

\end{document}
